%% =====================================================================
%%  T-12-03  Friends and Enemies
%%  -------------------------------------------------------------------
%%  One idea per file. Core is mandatory. Slots in canonical order.
%%  Max 700 words. No layout commands. No filler. New source material
%%  for this entry goes in sources/B3/T-12-03.tex -- never by
%%  rewriting this file (docs/RULES.md R0.1).
%%  =====================================================================
%% @kind: topic
%% @id: T-12-03
%% @title: Friends and Enemies
%% @subtitle: Never put too much trust in friends; learn to use enemies
%% @chapter: c12
%% @order: 30
%% @status: stable
%% @sources: B3
%% @updated: 2026-09-19

\topic{T-12-03}{Friends and Enemies}{Never put too much trust in friends; learn to use enemies}

\begin{core}
The counter-intuitive staffing rule: friendship is a poor qualification for the jobs that matter, and enmity is a workable one. Friends, mixed with power, drown in ingratitude and entitlement --- and you let your caution drop around them precisely because they are friends. A former enemy, reconciled, has everything to prove and no illusions to outgrow. Keep friends for friendship; hire for the job.
\end{core}
\begin{mechanism}
B3's mechanism is the jaws of ingratitude: you would never put a finger in a lion's mouth, but with friends every guard drops --- and hire one, and ingratitude eats you alive, because the friend you elevate stops seeing the gift and starts seeing the deserved. The psychology runs both ways: the friend expects treatment a stranger could not ask, resents discipline a stranger would accept, and envy --- the quiet engine --- works hardest at short distance (T-09-04). The reconciled enemy inverts every term: their reconciliation is on record, their loyalty is a performance they must keep staging to keep the reconciliation alive, and they bring their competence without the entitlement. The street proverb names the operating doctrine --- keep your friends where you can see them and your enemies where they can see you: closer, and legible. Absalom's revolt and the whole lineage of friend-betrays-king cases stand behind the law; Mao's use of enemies stands in front of it.
\end{mechanism}
\begin{conditions}
Strongest wherever power and staffing mix --- promotions, partnerships, family firms, governments. The friend-hire works only where accountability is explicit and the friend accepts subordination in writing. It weakens in genuinely flat partnerships among equals, and it does not apply to friendship itself --- the law is about mixing registers, not about people being untrustworthy.
\end{conditions}
\begin{application}
  \item Separate the ledgers: friendship is paid in time and honesty; work is paid in money and title. Never let one currency counterfeit the other.
  \item If you must work with friends, write the terms as if they were strangers --- salaries, exits, review dates. The paper protects the friendship, not the business.
  \item Before elevating a friend, simulate their first act of discipline from you. If you cannot administer it, do not appoint them.
  \item When you reconcile with an enemy, give them a job worth performing loyalty for --- and watch the performance like the theatre it is.
  \item Expect less from friends and you will not be bitter: never expecting gratitude from a friend means being pleasantly surprised when it comes.
\end{application}
\begin{feedback}
  \item Working: the friendship survived the working arrangement --- because the paper held and neither party tested it.
  \item Working: your reconciled enemy delivers, visibly over-performing to keep the record intact.
  \item Failing: the friend you promoted now treats direction as betrayal. The registers have fused.
  \item Failing: you find yourself unable to criticise, price, or fire someone for friendship's sake. They own the register now.
\end{feedback}
\begin{failure}
The doctrine calcifies into paranoia: a person who sees enemies everywhere recruits them by treatment, and the cynicism becomes self-fulfilling. It also has a loyalty cost on its own terms --- people notice who gets hired on merit and who on history, and the reconciled-enemy hire tells your loyalists that betrayal pays better than constancy. And friendship mixed with power can destroy the friendship in both directions even when the paper holds.
\end{failure}
\begin{countermeasures}
  \item Audit your own inner circle: who is there from history and who from merit? Neither is wrong --- mixing them unlabelled is what breaks organisations.
  \item When a rival sues for peace, accept it professionally and watch the record, not the words. Reconciliation is a contract paid in behaviour.
  \item Notice your caution dropping around friends --- the lion's jaw opens exactly there. Price friends' business proposals as you would a stranger's.
  \item If you are the friend being hired, insist on the stranger's paper yourself. It is the only way the friendship survives the job (T-04-02).
\end{countermeasures}

\sources{B3}
\seesources
\seealso{T-12-02, T-09-04, T-08-01}
